\documentclass[10pt,a3paper,landscape]{article}
\usepackage[margin=6mm]{geometry}
\usepackage[T1]{fontenc}
\usepackage[utf8]{inputenc}
\usepackage[french]{babel}
\usepackage{lmodern}
\usepackage{microtype}
\makeatletter\def\input@path{{../../../Style/}}\makeatother
\NeedsTeXFormat{LaTeX2e}
\ProvidesPackage{TLPosterCarteMentale}[2026/08/03 Posters et cartes mentales SI]

\RequirePackage{tikz}
\RequirePackage[most]{tcolorbox}
\RequirePackage{xcolor}
\RequirePackage{amsmath,amssymb}
\RequirePackage{graphicx}
\usetikzlibrary{positioning,arrows.meta,calc,shapes.geometric,mindmap,backgrounds,fit}

\definecolor{TLPosterBlue}{RGB}{24,86,141}
\definecolor{TLPosterLight}{RGB}{234,243,250}
\definecolor{TLPosterAccent}{RGB}{232,126,34}
\definecolor{TLPosterGreen}{RGB}{52,152,102}
\definecolor{TLPosterRed}{RGB}{192,57,43}
\definecolor{TLPosterGray}{RGB}{80,90,100}

\newcommand{\TLPosterTitre}[2]{%
  \begin{tcolorbox}[enhanced,boxrule=0pt,arc=0pt,colback=TLPosterBlue,left=6mm,right=6mm,top=4mm,bottom=4mm]
    {\color{white}\bfseries\LARGE #1}\hfill{\color{white}\large #2}
  \end{tcolorbox}%
}

\newtcolorbox{TLPosterBloc}[2][]{
  enhanced,breakable,
  colback=TLPosterLight,
  colframe=TLPosterBlue,
  colbacktitle=TLPosterBlue,
  coltitle=white,
  fonttitle=\bfseries,
  title={#2},
  boxrule=0.8pt,
  arc=2mm,
  left=3mm,right=3mm,top=2mm,bottom=2mm,
  #1
}

\newtcolorbox{TLPosterFormule}[1][]{
  enhanced,
  colback=white,
  colframe=TLPosterAccent,
  boxrule=1pt,
  arc=2mm,
  fontupper=\bfseries,
  halign=center,
  #1
}

\newcommand{\TLPosterMethode}[1]{%
  \begin{tcolorbox}[enhanced,colback=TLPosterGreen!8,colframe=TLPosterGreen,title=Méthode,fonttitle=\bfseries]
  #1
  \end{tcolorbox}%
}

\newcommand{\TLPosterAttention}[1]{%
  \begin{tcolorbox}[enhanced,colback=TLPosterRed!6,colframe=TLPosterRed,title=Point de vigilance,fonttitle=\bfseries]
  #1
  \end{tcolorbox}%
}

\tikzset{
  TLmindroot/.style={concept,concept color=TLPosterBlue,text=white,font=\bfseries\large,minimum size=34mm},
  TLmindlevel1/.style={level 1 concept/.append style={font=\bfseries,minimum size=23mm,sibling angle=45}},
  TLmindlevel2/.style={level 2 concept/.append style={font=\small,minimum size=17mm}},
  TLarrow/.style={-{Stealth[length=3mm]},thick,draw=TLPosterBlue},
  TLbox/.style={rounded corners=2mm,draw=TLPosterBlue,fill=TLPosterLight,align=center,inner sep=3mm}
}

\newenvironment{TLCarteMentale}[1]{%
  \begin{tikzpicture}[mindmap,grow cyclic,concept color=TLPosterBlue,every node/.style={concept},TLmindlevel1,TLmindlevel2]
  \node[TLmindroot]{#1}
}{;
  \end{tikzpicture}
}

\newcommand{\TLBranche}[3][]{child[concept color=#2]{node[#1]{#3}}}

\endinput

\pagestyle{empty}
\renewcommand{\familydefault}{\sfdefault}
\begin{document}
\begin{tikzpicture}
\TLMapBase{Cinématique}{Position $\rightarrow$ Vitesse $\rightarrow$ Accélération}
\TLMapBranchTL{Savoirs}{\textbullet\ maîtriser produit vectoriel, produit mixte et repérage spatial\\[1.2mm]\textbullet\ définir un torseur cinématique et effectuer un changement de point\\[1.2mm]\textbullet\ caractériser translation, rotation, mouvement plan et trajectoires\\[1.2mm]\textbullet\ relier position, vitesse, accélération et dérivation vectorielle\\[1.2mm]\textbullet\ composer les mouvements et traiter contact, roulement sans glissement et profil trapézoïdal}
\TLMapBranchTR{Outils}{\textbullet\ produit vectoriel\\[1.2mm]\textbullet\ produit mixte\\[1.2mm]\textbullet\ repères}
\TLMapBranchML{Torseur cinématique}{\textbullet\ rotation\\[1.2mm]\textbullet\ vitesse en un point\\[1.2mm]\textbullet\ changement de point}
\TLMapBranchMR{Grandeurs}{\textbullet\ position\\[1.2mm]\textbullet\ vitesse\\[1.2mm]\textbullet\ accélération\\[1.2mm]\textbullet\ dérivation vectorielle}
\TLMapBranchBL{Applications}{\textbullet\ composition\\[1.2mm]\textbullet\ roulement sans glissement\\[1.2mm]\textbullet\ loi trapézoïdale}
\TLMapBranchBR{Méthode}{1. Choisir solides, repères, points et paramètres adaptés.\\[1.2mm]2. Écrire les vecteurs position et les relations géométriques utiles.\\[1.2mm]3. Dériver dans le bon repère pour obtenir vitesse puis accélération.\\[1.2mm]4. Transporter les vitesses au point pertinent et composer les mouvements.\\[1.2mm]5. Appliquer les conditions de contact ou de roulement sans glissement.}
\TLMapRibbon{Repérer $\rightarrow$ Exprimer $\rightarrow$ Dériver $\rightarrow$ Composer $\rightarrow$ Interpréter}
\end{tikzpicture}
\end{document}
